\documentclass[10pt,aspectratio=169]{beamer}

\usepackage[utf8]{inputenc}
\usepackage[T1]{fontenc}
\usepackage[french]{babel}
\usepackage{booktabs}
\usepackage{tabularx}
\usepackage{hyperref}

\usetheme{Madrid}
\usecolortheme{default}
\setbeamertemplate{navigation symbols}{}

\definecolor{marketgreen}{RGB}{33,95,70}
\definecolor{marketlight}{RGB}{233,242,237}
\definecolor{marketaccent}{RGB}{190,121,44}

\setbeamercolor{palette primary}{bg=marketgreen,fg=white}
\setbeamercolor{palette secondary}{bg=marketgreen,fg=white}
\setbeamercolor{palette tertiary}{bg=marketgreen,fg=white}
\setbeamercolor{structure}{fg=marketgreen}
\setbeamercolor{frametitle}{bg=marketlight,fg=marketgreen}
\setbeamercolor{block title}{bg=marketgreen,fg=white}
\setbeamercolor{block body}{bg=marketlight,fg=black}

\title{Sprint 1 Review}
\subtitle{Marketplace Locale Intelligente}
\author{Thomas\\Lead Developer et Product Owner}
\date{Mars 2026}

\begin{document}

\begin{frame}
  \titlepage
\end{frame}

\begin{frame}{Objectif du sprint 1}
  \begin{block}{Ambition}
    Poser un socle produit et technique suffisamment propre pour lancer la logique métier au sprint 2.
  \end{block}

  \vspace{0.4cm}

  \textbf{Périmètre visé}
  \begin{itemize}
    \item inscription client et producteur;
    \item authentification et récupération du profil;
    \item consultation du catalogue produit;
    \item mise en place des fondations backend, frontend, DevOps et data.
  \end{itemize}
\end{frame}

\begin{frame}{Backlog ciblé}
  \begin{columns}[T]
    \begin{column}{0.48\textwidth}
      \textbf{User stories}
      \begin{itemize}
        \item \textbf{US1.1} Créer un compte client ou producteur
        \item \textbf{US1.2} Se connecter et récupérer son profil
        \item \textbf{US1.3} Consulter les produits et leur détail
      \end{itemize}
    \end{column}
    \begin{column}{0.48\textwidth}
      \textbf{Features techniques}
      \begin{itemize}
        \item backend + PostgreSQL
        \item auth + CRUD produits
        \item frontend React/Vite
        \item GitFlow + CI initiale
        \item Docker + compose
        \item mock data pour la data team
      \end{itemize}
    \end{column}
  \end{columns}
\end{frame}

\begin{frame}{Livrables backend et base}
  \textbf{Fondations livrées}
  \begin{itemize}
    \item API FastAPI structurée en routeurs versionnés \texttt{/api/v1}
    \item modèles SQLAlchemy et configuration base de données
    \item schéma PostgreSQL versionné
    \item dépendances auth et schémas API
  \end{itemize}

  \vspace{0.3cm}

  \textbf{Endpoints démontrables}
  \begin{itemize}
    \item \texttt{POST /auth/register}
    \item \texttt{POST /auth/login}
    \item \texttt{GET /users/me}
    \item \texttt{GET /products}
    \item \texttt{GET /products/\{id\}}
    \item \texttt{GET /producers}
  \end{itemize}
\end{frame}

\begin{frame}{Livrables frontend}
  \textbf{Socle UI livré}
  \begin{itemize}
    \item application React/Vite structurée par pages, composants et contextes
    \item pages d'accueil, login, register, catalogue et producteurs
    \item consommation des endpoints d'authentification et de catalogue
    \item base prête pour le tunnel panier / commande du sprint 2
  \end{itemize}

  \vspace{0.3cm}

  \begin{block}{Constat PO}
    Le frontend n'est pas seulement initialisé: il fournit déjà une base démontrable et directement exploitable pour la suite.
  \end{block}
\end{frame}

\begin{frame}{Livrables DevOps et data}
  \begin{columns}[T]
    \begin{column}{0.48\textwidth}
      \textbf{DevOps}
      \begin{itemize}
        \item Dockerfile backend
        \item Dockerfile frontend multi-stage
        \item \texttt{docker-compose} front / back / postgres
        \item workflow CI backend
        \item règles GitFlow documentées
      \end{itemize}
    \end{column}
    \begin{column}{0.48\textwidth}
      \textbf{Data}
      \begin{itemize}
        \item schéma pensé pour le transactionnel et l'analytique
        \item script de seed volumineux
        \item notebook d'exploration initial
        \item documentation dédiée au dataset
      \end{itemize}
    \end{column}
  \end{columns}
\end{frame}

\begin{frame}{Bilan par user story et feature}
  \small
  \begin{tabularx}{\textwidth}{>{\bfseries}l l X}
    \toprule
    ID & Statut & Commentaire \\
    \midrule
    US1.1 & Livrée & Inscription disponible côté API et frontend \\
    US1.2 & Livrée & Login JWT et profil utilisateur opérationnels \\
    US1.3 & Partielle & Liste produits OK, détail produit incomplet côté frontend \\
    FEAT1.1 & Livrée & Backend et schéma PostgreSQL en place \\
    FEAT1.2 & Livrée & Auth, profil, produits, producteurs disponibles \\
    FEAT1.3 & Livrée & Frontend structuré et exploitable \\
    FEAT1.4 & Partielle & GitFlow et CI présents, couverture qualité encore limitée \\
    FEAT1.5 & Livrée & Dockerfiles et environnement compose disponibles \\
    FEAT1.6 & Livrée & Mock data et support exploration data présents \\
    \bottomrule
  \end{tabularx}
\end{frame}

\begin{frame}{Points forts du sprint}
  \begin{itemize}
    \item fondations critiques intégrées dès le premier sprint;
    \item bonne cohérence d'ensemble entre application, infra et data;
    \item architecture déjà préparée pour absorber le sprint 2;
    \item base crédible pour une démonstration produit et technique.
  \end{itemize}

  \vspace{0.3cm}

  \begin{block}{Lecture PO}
    Le sprint 1 a bien rempli son rôle: réduire le risque projet en installant un socle transverse, pas seulement une collection de fichiers initiaux.
  \end{block}
\end{frame}

\begin{frame}{Écarts et risques observés}
  \begin{itemize}
    \item détail produit non finalisé côté expérience frontend;
    \item couverture de tests encore faible;
    \item CI centrée sur le backend, pas encore équilibrée avec le frontend;
    \item quelques routes frontend préparées avant finalisation complète des écrans.
  \end{itemize}

  \vspace{0.3cm}

  \textbf{Impact}
  \begin{itemize}
    \item risque de décalage entre “présent dans le code” et “démontrable sans réserve”;
    \item besoin de consolidation rapide avant d'ouvrir davantage le périmètre métier.
  \end{itemize}
\end{frame}

\begin{frame}{Décision Product Owner}
  \begin{block}{Décision}
    Sprint 1 globalement accepté.
  \end{block}

  \textbf{Accepté}
  \begin{itemize}
    \item socle backend, base, Docker et data;
    \item socle frontend et intégration initiale;
    \item gouvernance Git et première CI.
  \end{itemize}

  \textbf{Acceptation partielle}
  \begin{itemize}
    \item complétude de \textbf{US1.3};
    \item niveau de validation automatisée attendu pour la suite.
  \end{itemize}
\end{frame}

\begin{frame}{Priorités pour le sprint 2}
  \begin{enumerate}
    \item finaliser les écarts de finition hérités du sprint 1;
    \item sécuriser le flux catalogue $\rightarrow$ panier $\rightarrow$ commande;
    \item renforcer la CI avec build frontend et plus de tests;
    \item stabiliser l'intégration front-back avant d'élargir le périmètre.
  \end{enumerate}

  \vspace{0.3cm}

  \begin{block}{Point de vigilance}
    Le prochain sprint doit transformer un socle prometteur en parcours métier cohérents de bout en bout.
  \end{block}
\end{frame}

\begin{frame}{Conclusion}
  \begin{itemize}
    \item sprint 1 réussi sur la mise en place des fondations;
    \item architecture, frontend, DevOps et data désormais raccordés;
    \item quelques écarts de finition à fermer rapidement;
    \item base saine pour attaquer la logique métier du sprint 2.
  \end{itemize}

  \vspace{0.5cm}

  \centering
  \Large Questions / Validation PO
\end{frame}

\end{document}